\section{방법론}
\label{sec:method}

% P0: 그림 1의 구조와 절의 대응
\rev{그림~\ref{fig:pipeline}은 본 연구의 절차를 네 블록으로 보인다.
데이터 수집(A)과 누수 없는 표본 구성(B)은 공개 텔레메트리를 팀--단계 예측 표본으로 바꾸고, 설계와 분석(C)은 세 연구 질문에 대응하는 모델을 구성하며, 평가(D)는 모든 설계를 같은 조건에서 비교한다.
이하의 각 절은 그림의 블록과 같은 순서로 서술한다.}

\begin{figure*}[t]
\centering
\includegraphics[width=\textwidth]{fig_pipeline_clean.pdf}
\caption{데이터 파이프라인과 모델링·평가 절차}
\label{fig:pipeline}
\end{figure*}

%======================================================================
\subsection{데이터 수집}
\label{sec:method:data}

% P1: 무엇이 기록되는가
PUBG는 공개 API를 통해 경기마다 텔레메트리를 제공한다~\cite{PUBGAPI}.
텔레메트리에는 경기에 참가한 모든 플레이어의 위치가 약 10초 간격으로 기록되며, 피해, 기절, 처치, 아이템의 획득·장착·사용, 안전구역의 상태 변화, 단계 전환, 착지와 같은 사건이 발생 시각과 함께 기록된다.
따라서 하나의 경기 안에서 모든 팀의 이동 궤적과 팀 사이의 교전, 그리고 안전구역의 변화를 함께 관측할 수 있다.

% P2: 수집 방식, 규모, 희소성
본 연구는 \rev{공개 API의 표본 경기 목록을 하루 단위로 조회하여} 스팀 서버의 경기를 수집하였다.
\rev{수집한 경기는 모두 2026년 6월의 안전구역 규칙 개편~\cite{PUBGBlueZone2026} 이전에 생성되어, 모두 개편 전의 같은 안전구역 규칙을 따른다.}
\rev{분석에 사용한 기록의 유형별 규모는 표~\ref{tab:corpus}에, 코호트별 경기와 표본의 수는 표~\ref{tab:cohorts}에 정리하였다.}
경기 종료 후의 통계를 집계한 공개 데이터셋은 존재하지만~\cite{KagglePUBG}, 경기 중 모든 플레이어의 위치를 이 해상도로 담은 자료는 드물며, 기존 배틀로얄 예측 연구가 사용한 경기 내 기록도 공개되지 않았다~\cite{Sagong25a,KRAFTON}.

\begin{table}[t]
\centering
\caption{분석에 사용한 경기의 텔레메트리 기록 유형별 건수}
\label{tab:corpus}
\footnotesize
\begin{tabular}{lr}
\hline
기록 유형 & 건수 \\
\hline
위치 표본 (약 10초 간격) & 96,426,398 \\
피해 & 90,178,937 \\
아이템 획득·버림 & 67,865,634 \\
아이템 장착 & 46,218,920 \\
아이템 사용 & 15,229,394 \\
안전구역 상태 갱신 & 2,134,448 \\
착지 & 1,673,902 \\
처치 & 1,267,179 \\
기절 & 1,199,290 \\
단계 전환 & 200,979 \\
\hline
합계 & 322,395,081 \\
\hline
\end{tabular}
\end{table}

%======================================================================
\subsection{누수 없는 표본 구성}
\label{sec:method:samples}

\rev{원시 기록을 팀--단계 예측 표본으로 바꾸는 과정은 그림~\ref{fig:pipeline}의 B와 같이 시간축 정렬, 예측 시점과 입력 절단, 사건 기반 라벨, 표본 구성의 순서로 이루어진다.}

\subsubsection{공통 시간축과 단계 경계}
위치 기록에 함께 저장된 경기 경과 시간은 사건 기록의 타임스탬프와 기준시각이 달라 경기마다 어긋난다.
본 연구는 모든 기록의 타임스탬프를 경기 시작 사건으로부터의 경과 시간으로 변환하여 하나의 시간축을 사용한다.
단계 경계는 이 시간축 위의 단계 전환 사건으로 정한다.
텔레메트리는 새 안전구역의 공개와 구역 축소의 시작을 단계 전환 사건으로 기록하며, 한 구역으로의 축소가 끝나는 시각에 다음 구역이 공개된다~\cite{PUBGAPI}.
본 연구는 $k$번째 안전구역이 공개된 시점부터 그 구역으로의 축소가 끝날 때까지를 단계 $k$로 정의하고, 축소가 끝나 다음 안전구역이 공개되는 시각을 단계 $k$의 종료 시각 $b_k$로 둔다.

\subsubsection{예측 시점과 입력 절단}
예측 시점은 단계 경계 $b_k$이며, 모든 입력은 그보다 10초 앞선 특징 절단 시각 $a_k=b_k-10$초 이전의 기록만으로 구성한다.
이 제한은 팀 특징뿐 아니라 생존 인원, 팀 사이의 관계, attention 마스크에도 적용된다.
생존 여부 역시 최종 탈락 시각이 아니라 $a_k$ 이전의 처치와 위치 기록으로 판정하며, 처치된 플레이어는 처치 후 20초 이상 지난 시점에 체력이 남은 위치 기록이 나타날 때에만 재투입된 것으로 본다.
\rev{입력이 절단 시각 이후의 기록에 의존하지 않음을 확인하기 위해, 무작위로 고른 경기에서 $a_k$ 이후의 원시 기록을 모두 지우고 입력을 다시 구성하여 원래 입력과 비교하였으며, 모든 경계에서 두 입력이 일치하였다.}

\subsubsection{사건 기반 탈락 라벨}
탈락 라벨은 사건 기록으로 정의한다.
기절은 이후 소생이 가능한 상태이므로 최종 처치와 구분하며, 팀의 탈락 시각 $T$는 재투입되지 않은 마지막 팀원이 처치된 시각으로 정의한다.
이렇게 정한 탈락 시각은 게임이 팀의 최종 순위를 기록한 처치 사건의 시각과 일치하였다.
우승 팀은 탈락 사건 없이 경기가 끝나므로 경기 종료 시각 $C$에서 관측을 종료한다.
단계 $k$의 위험집합은 $b_k$에 생존한 팀들이며, 각 팀의 라벨은 다음 단계의 종료 시각 $b_{k+1}$까지 탈락하면 1, 그렇지 않으면 0이다.
다음 단계가 경기의 마지막 단계이면 $b_{k+1}$ 대신 경기 종료 시각을 사용한다.
\rev{경기 종료 기록이 없거나 라벨 구간에서 위치 기록이 60초 넘게 끊긴 경기는 제외하도록 하였으나, 해당하는 경기는 없었다.}

\subsubsection{팀--단계 표본과 예측 대상}
\rev{예측 시점은 단계 1부터 7까지의 경계에 두며, 분석의 단위는 경기 $m$, 팀 $i$, 단계 경계 $k$의 조합인 팀--단계 표본이다.
팀은 1인, 2인 또는 4인으로 구성되므로, 솔로 경기에서는 팀--단계 표본이 곧 플레이어의 단계별 표본이다.}
예측 대상은 다음 단계 탈락 위험
\begin{equation}
\lambda_{m,i,k}=\Pr\!\left(T_{m,i}\le b_{m,k+1} \mid T_{m,i}>b_{m,k},\,\mathcal{H}_{m}(a_{m,k})\right)
\label{eq:hazard}
\end{equation}
이며, $\mathcal{H}_{m}(t)$는 시각 $t$까지의 경기 기록이다.
이 정의는 각 단계 경계를 landmark로, 다음 단계를 예측 구간으로 하는 이산시간 hazard에 해당한다~\cite{SingerWillett93,vanHouwelingen07}.
위험집합의 구성에는 예측 시점의 생존 여부를 사용하지만, 모델의 입력은 모두 특징 절단 시각 이전의 정보로만 구성한다.

\subsubsection{코호트와 분할}
\label{sec:method:cohort}
분석 대상은 일반 경기와 랭크 경기로 한정하였다.
실제 경쟁이 이루어지지 않는 튜토리얼과 훈련장 경기는 제외하였고, 일반적인 경기 양상과 다른 두 모집단, 즉 \rev{계정 대부분이 AI인} airoyale 경기와 \rev{비플레이어 전투 계정이 다수 기록되는} Vikendi 맵의 경기도 제외하였다.
일반 경기에도 AI 계정이 섞여 있으므로, 각 팀의 AI 계정 비율을 입력 특징에 포함하였다.
경기는 맵, 시점(3인칭·1인칭), 팀 규모(1인·2인·4인)의 조합으로 코호트를 나누었으며, airoyale 경기를 제외하기 전 기준으로 300경기 이상인 조합을 코호트로 정하였다\rev{(표~\ref{tab:cohorts})}.
각 코호트의 경기는 경기 단위로 학습, 검증, 시험 집합에 8:1:1로 나누었고, 분할은 고정하여 모든 비교에 같은 분할을 사용한다.

\begin{table}[t]
\centering
\caption{\rev{코호트별 경기 수, 팀--단계 표본 수, 다음 단계 탈락률}}
\label{tab:cohorts}
\footnotesize
\setlength{\tabcolsep}{3pt}
\begin{tabular}{lccrrrr}
\hline
맵 & 시점 & 인원 & 경기 & 학습/검증/시험 & 표본 & 탈락률(\%) \\
\hline
Erangel & 3인칭 & 1 & 224 & 189/15/20 & 31,740 & 35.1 \\
Erangel & 3인칭 & 2 & 656 & 526/65/65 & 59,170 & 29.5 \\
Taego & 3인칭 & 2 & 573 & 459/57/57 & 47,223 & 29.2 \\
Erangel & 1인칭 & 4 & 430 & 344/43/43 & 21,849 & 27.8 \\
Erangel & 3인칭 & 4 & 3,283 & 2,638/321/324 & 196,442 & 25.7 \\
Karakin & 3인칭 & 4 & 333 & 267/33/33 & 12,669 & 32.1 \\
Miramar & 3인칭 & 4 & 1,543 & 1,235/154/154 & 92,279 & 24.6 \\
Paramo & 3인칭 & 4 & 333 & 267/33/33 & 9,963 & 30.5 \\
Rondo & 3인칭 & 4 & 1,039 & 833/103/103 & 53,683 & 25.2 \\
Sanhok & 3인칭 & 4 & 940 & 752/94/94 & 48,793 & 27.6 \\
Taego & 1인칭 & 4 & 435 & 349/43/43 & 23,390 & 30.0 \\
Taego & 3인칭 & 4 & 4,159 & 3,329/415/415 & 220,565 & 24.9 \\
\hline
합계 & & & 13,948 & 11,188/1,376/1,384 & 817,766 & 26.6 \\
\hline
\end{tabular}
\\[2pt]
\parbox{\columnwidth}{\scriptsize 2026년 3월 26일--6월 2일에 생성된 경기. 표본은 위험집합에 속한 팀--단계 표본이며, 탈락률은 그중 다음 단계 안에 탈락한 비율이다.}
\end{table}

%======================================================================
\subsection{설계와 분석}
\label{sec:method:design}

\rev{그림~\ref{fig:pipeline}의 C는 세 입력 표현과 세 연구 질문에 대응하는 분석을 보인다.
서식지 맥락 표현은 RQ1의 해석 모델과 RQ2의 본 연구 설계에 입력되고, 두 기존 설계의 입력은 RQ2의 적용판과 RQ3의 분해에 사용된다.}

\subsubsection{서식지 맥락 표현}
\label{sec:method:representation}
서식지 축소 관점에서 각 팀의 상태를 네 가지 구성개념으로 정의한다.
첫째, 구역 경계에 대한 노출은 팀 중심과 안전구역 경계 사이의 부호 있는 거리로 측정하며, 구역 밖에 있는 깊이와 구역 안쪽으로 들어온 깊이로 나누어 사용한다.
둘째, 이동 부담은 안전구역에 도달하는 데 필요한 속도와 팀이 실제로 이동한 속도의 차이, 즉 이동 부족분으로 측정한다~\cite{Nathan08}.
필요 속도는 특징 절단 시각 직전의 팀 중심에서 그 시점에 공개된 안전구역 경계까지의 거리(구역 안이면 0)를 단계 길이로 나눈 값이고, 이동 속도는 단계 시작부터 특징 절단 시각까지의 이동 경로 길이를 같은 단계 길이로 나눈 값이다.
단계 길이에는 경기별 관측값 대신 학습 집합에서 코호트, 경기 유형, 단계별로 구한 중앙값을 사용하여, 특징 절단 시각 이후의 정보가 들어가지 않게 하였다.
셋째, 주변 경쟁 밀도는 팀 주변 200\,m 안에 있는 다른 팀의 수이다.
넷째, 남은 팀원 수는 특징 절단 시각에 생존한 것으로 관측된 팀원의 수이다.
\rev{필요 속도는 구역 밖 깊이와 상관이 높아 정보가 중복되므로, 실제 이동을 함께 반영하는 이동 부족분을 사용하였으며}, 이 선택은 결과를 확인한 뒤 이루어진 탐색적 특징 설계이다.

예측 모델의 입력은 이 구성개념을 포함한 팀--단계 특징으로, 경기 맥락, 팀 상태, 안전구역, 이동, 노출, 자원, 공간, 관계, 서식지의 아홉 계열과 팀 중심 좌표, 팀의 AI 계정 비율로 이루어진다.
같은 단계의 팀들 사이에는 세 유형의 관계를 정의한다.
근접 관계는 단계 시작부터 특징 절단 시각까지 두 팀의 평균 거리 $d$에 대해 가중치 $\exp(-d/200)$를 갖는 대칭 관계이다.
교전 관계는 같은 구간에서 한 팀이 다른 팀에 가한 피해, 기절, 처치를 결합한 강도 $\log(1+\text{강도})$를 가중치로 갖는 방향 관계이며, 교전이 있었던 팀 쌍에서는 근접 관계를 교전 관계로 대체한다.
공동 압박 관계는 두 팀이 모두 구역 밖에 있거나 경계 안쪽 150\,m 이내에 있으면서 서로 200\,m 이내에 있을 때의 대칭 관계이다.

\subsubsection{RQ1: 생태 hazard 모델}
\label{sec:method:rq1}
RQ1에 답하기 위해 서식지 맥락의 구성개념만으로 위험을 설명하는 로지스틱 회귀를 추정한다.
\begin{equation}
\mathrm{logit}\,\lambda_{m,i,k}=\alpha_k+\sum_{c}\beta_c z_{c,m,i,k}
\label{eq:ehm}
\end{equation}
여기서 $\alpha_k$는 단계별 기저위험이고, $z_c$는 학습 집합에서 표준화한 구성개념(구역 밖 깊이, 구역 안쪽 깊이, 이동 부족분, 주변 경쟁 밀도, 남은 팀원 수)이며, 전체 생존 팀 수의 로그를 통제 변수로 포함한다.
솔로 코호트에서는 남은 팀원 수가 상수이므로 제외한다.
모델은 코호트마다 학습 집합에서 추정하며, 표준오차는 경기 단위로 군집화하여 계산한다.
코호트별 계수는 확률효과 모형으로 통합하고~\cite{DerSimonian86}, 통합 추정치와 95\% 신뢰구간, 코호트 간 이질성, 그리고 계수의 부호가 통합 추정치와 같은 코호트의 수를 보고한다.
단계에 따른 변화를 보기 위해 각 구성개념과 단계 구간(초반 1--2, 중반 3--5, 후반 6--7)의 상호작용을 추가한 모델을 추정하고, 구간 간 기울기가 같은지 검정한다.

\subsubsection{RQ2: 단계 조건부 hazard 설계와 기존 설계}
\label{sec:method:rq2}
본 연구의 설계는 식~\eqref{eq:hazard}의 위험을 로짓 척도에서 단계별 기저위험과 팀 상태의 함수
\begin{equation}
\mathrm{logit}\,\lambda_{m,i,k}=\alpha_k+f(\mathbf{x}_{m,i,k},\mathcal{G}_{m,k})
\label{eq:logit}
\end{equation}
로 두고, 다음 단계 탈락에 대한 교차엔트로피로 학습한다.
여기서 $\mathbf{x}_{m,i,k}$는 팀의 서식지 맥락, $\mathcal{G}_{m,k}$는 같은 단계의 다른 팀들과의 관계이다.
설계가 특정 학습기에 의존하지 않음을 보이기 위해 두 학습기로 구현한다.
첫째는 \rev{그래디언트 부스팅 결정나무(GBDT)}~\cite{Friedman01,Ke17}로, 단계 번호를 입력 특징으로 포함하며 $\alpha_k$를 별도의 항으로 두지는 않는다.
둘째는 \rev{구조 요소의 기여를 분해하기 위한} 참조 모델인 PTR-Former로, 팀--단계 토큰에 단계 임베딩을 더한 뒤 같은 팀의 단계 이력에 대한 인과적 attention(단계 축)과 같은 단계의 생존 팀들에 대한 attention(팀 축)을 번갈아 적용하고~\cite{Vaswani17}, 팀 축 attention의 로짓에는 관계 유형별로 학습되는 편향을 더한다~\cite{Ying21}.
PTR-Former는 hazard의 교차엔트로피에 최종 순위에 대한 쌍별 순위 손실을 작은 가중치로 더하여 학습한다\footnote{차원 128, attention 헤드 4개, 블록 2개, dropout 0.1, 순위 손실 가중치 0.1을 사용하였다. AdamW~\cite{LoshchilovAdamW}와 코사인 학습률 감소~\cite{LoshchilovSGDR}로 최대 12 epoch 동안 학습하였고, 검증 집합의 로그 손실이 3 epoch 동안 개선되지 않으면 학습을 멈추었다.}.

\rev{기존 설계 가운데 KRAFTON의 예측기는 같은 위험집합 안에서 팀들의 점수 차이를 학습하는 상대위험 목적을 사용한다~\cite{KRAFTON}.}
Cox 부분우도의 각 항은 사건 팀의 점수가 위험집합에서 차지하는 비율 $\exp(s_i)/\sum_{j}\exp(s_j)$이므로~\cite{Cox72}, 이 목적은 같은 시점의 모든 점수에 공통으로 더해지는 이동에 불변이며 그 자체로는 식~\eqref{eq:logit}의 $\alpha_k$와 같은 단계별 공통 위험 수준을 식별하지 못한다.
기저위험을 따로 추정하여 결합하면 구간별 확률을 얻을 수 있으므로~\cite{PrenticeGloeckler78}, 상대위험 점수에 단계별 기저위험을 결합한 변형도 함께 비교한다.

두 기존 설계는 원 설계의 입력과 학습 목적을 유지하면서 본 연구의 표본에 맞게 적용하였으며, \rev{두 설계와 본 연구 설계의 차이는 표~\ref{tab:compare}에 정리하였다.}
Sagong 등의 설계~\cite{Sagong25a}는 플레이어 단위 지표로 구성되므로 2인·4인 코호트에서는 팀 단위로 집계하였고\rev{, 솔로 코호트에서는 원 설계와 같은 개인 단위를 그대로 사용한다}.
정규화 자료에 없는 공격 시도와 어시스트 기록은 가장 가까운 사건으로 대체하였다.
KRAFTON의 예측기는 공개 구현을 바탕으로 입력 지표를 본 연구의 텔레메트리로 재구성하였고, \rev{출력은 검증 집합에서 확률로 보정하였다~\cite{Platt99}.}
모든 설계는 같은 경기 분할을 사용하고 같은 팀--단계 위험집합과 라벨에서 평가하며, RQ2에서는 각 적용판의 원래 학습 표집과 학습 목적을 유지한다.

적용판이 본 연구의 과제에 불리하게 구성되지 않았는지 확인하기 위해, 각 적용판을 원래 과제로도 평가한다.
KRAFTON 적용판은 원 구현이 정의한 우승 팀 정확도와 탈락 순서의 일치도(C-index)로, Sagong 적용판은 \rev{분석 단위가 같은 솔로 코호트에서} 개인 생존 분류로 평가한다.
본 연구의 설계도 같은 지표로 평가하여, 과제에 따라 설계 간 우열이 어떻게 달라지는지 확인한다.

\begin{table*}[t]
\centering
\caption{\rev{기존 설계와 본 연구 설계의 비교}}
\label{tab:compare}
\footnotesize
\renewcommand{\arraystretch}{1.2}
\begin{tabular}{p{2.3cm}p{4.3cm}p{4.6cm}p{4.6cm}}
\hline
 & Sagong 등~\cite{Sagong25a} & KRAFTON~\cite{KRAFTON} & 본 연구 \\
\hline
분석 단위 & 플레이어 & 팀 & 팀(솔로 경기에서는 개인) \\
예측 대상 & 다음 반 단계의 생존 여부 & 탈락 순서의 상대위험(우승확률로 변환) & 다음 단계 안에 탈락할 조건부 확률 \\
예측 시점 & 반 단계 경계 & 경기 중 여러 시점 & 단계 경계(다음 안전구역 공개 시점) \\
단계별 위험 수준 & 플레이 유형별 전이 생존율 & 표현하지 않음 & 단계별 기저위험으로 명시 \\
입력 & 단계별 활동·교전 점수, 플레이 유형 & 팀의 행동·상태 지표, 팀원 위치 & 서식지 맥락, 팀 간 관계 \\
팀 간 상호작용 & 다루지 않음 & 팀 간 attention & 팀 축·단계 축 attention과 관계 편향 \\
학습 목적 & 생존 분류 & 가중 Cox 부분우도 & 다음 단계 탈락의 교차엔트로피 \\
\hline
본 연구의 적용 & 2·4인 팀은 팀 단위로 집계, 예측 시점을 단계 경계로 통일 & 입력을 본 연구의 텔레메트리로 재구성, 출력을 확률로 보정 & --- \\
\hline
\end{tabular}
\end{table*}

\subsubsection{RQ3: 단일 요인 분해}
\label{sec:method:decomposition}
RQ3에 답하기 위해 다른 조건을 고정한 채 한 번에 하나의 설계 요인만 바꾸는 비교를 구성한다(표~\ref{tab:ladder}).
입력 표현의 효과는 학습기를 GBDT로 고정하고 기존 설계의 입력을 본 연구의 입력으로 바꾸어 측정하며, Sagong 등의 입력과의 비교는 서식지 맥락을 먼저 더한 뒤 나머지 입력을 바꾸는 두 단계로 나누어 서식지 맥락이 차지하는 부분을 구분한다.
\rev{학습 목적과 단계별 기저위험의 효과는 KRAFTON의 구조에 본 연구의 입력을 넣고 Cox 목적으로 학습한 모델에서 출발하여, 단계별 기저위험의 결합, 다음 단계 탈락의 교차엔트로피로의 학습 목적 교체, 명시적인 단계 절편의 추가를 차례로 비교한다.}
두 학습 목적의 모델은 구조, 입력, 학습 경기와 절단 시점, 초기 가중치, 갱신 횟수가 같으며, Cox 목적은 절단 시각에 생존한 팀들의 이후 탈락 순서를, 교차엔트로피 목적은 위험집합의 다음 단계 탈락 여부를 감독 정보로 사용한다.
KRAFTON 원본에서 이 출발 모델로 넘어가는 비교는 입력과 학습 표본이 함께 바뀌므로 분해 단계와 구분하여 보고한다.
\rev{이어서 같은 학습 목적과 입력에서 KRAFTON 구조를 팀 축만 남긴 PTR-Former로 바꾸어 아키텍처의 효과를, 여기에 단계 축과 관계 편향을 더하여 구조의 효과를 측정한다.}
이 두 비교에서는 학습 목적을 맞추기 위해 순위 보조 손실을 사용하지 않으며, 순위 보조 손실의 효과는 따로 보고한다.

구조 요인은 단계 축과 팀 축을 2$\times$2로 켜고 끄는 요인 실험으로도 평가하고, 관계의 대응을 경기 안에서 무작위로 섞어 다시 학습한 대조 모델로 관계 편향의 기여를 확인한다.
또한 안전구역의 원시 기하만을 갖는 구역 객체를 토큰으로 더하여 서식지 특징 없이도 그 정보를 학습하는지 확인하고, GBDT의 입력에 재구성한 KRAFTON 지표를 더하여 두 표현의 보완성을 확인한다.

\begin{table}[t]
\centering
\caption{\rev{단일 요인 분해의 단계}}
\label{tab:ladder}
\footnotesize
\setlength{\tabcolsep}{3pt}
\begin{tabular}{p{0.6cm}p{4.9cm}p{1.9cm}}
\hline
단계 & 비교 & 바꾸는 요인 \\
\hline
(a1) & GBDT: Sagong 입력 $\rightarrow$ 서식지 맥락 추가 & 서식지 맥락 \\
(a2) & GBDT: 서식지 맥락을 더한 Sagong 입력 $\rightarrow$ 본 연구 입력 & 나머지 입력 \\
(b) & GBDT: KRAFTON 입력 $\rightarrow$ 본 연구 입력 & 입력 표현 \\
(c1) & KRAFTON 구조·Cox 목적 $\rightarrow$ 단계별 기저위험 결합 & 기저위험 \\
(c2) & 기저위험을 결합한 Cox 목적 $\rightarrow$ 다음 단계 교차엔트로피 & 학습 목적 \\
(d) & 다음 단계 교차엔트로피 $\rightarrow$ 명시적 단계 절편 추가 & 단계 절편 \\
(e) & KRAFTON 구조 $\rightarrow$ 팀 축 구조 & 아키텍처 \\
(f) & 팀 축 구조 $\rightarrow$ 단계 축·관계 편향 추가(PTR-Former) & 구조 \\
(g) & GBDT: 본 연구 입력 $\rightarrow$ KRAFTON 입력 추가 & 입력 확장 \\
\hline
\end{tabular}
\\[2pt]
\parbox{\columnwidth}{\scriptsize 각 단계는 표시한 요인만 바꾸고, 학습기, 입력, 표본, 학습 목적 가운데 나머지는 직전 단계와 같다. (e)와 (f)는 순위 보조 손실 없이 비교한다.}
\end{table}

%======================================================================
\subsection{평가}
\label{sec:method:evaluation}

\rev{그림~\ref{fig:pipeline}의 D와 같이, 모든 설계의 예측은 같은 시험 표본에서 지표로 요약하고 코호트별 차이를 통합한다.}

\subsubsection{평가 지표}
주 지표는 시험 집합의 위험집합 전체에 대한 AUC이다.
단계 안 AUC는 같은 단계 번호의 표본을 경기에 걸쳐 모아 단계별로 계산한 AUC의 평균으로, 단계마다 공통으로 더해지는 위험 수준의 차이를 제외한 판별력을 나타내며, 시험 표본이 30개 이상이고 두 라벨이 모두 있는 단계만으로 계산한다.
경기--단계 내부 AUC는 같은 경기의 같은 단계 위험집합 안에서 팀을 구분하는 판별력으로, 위험집합별 AUC를 양성--음성 쌍의 수로 가중하여 통합하고 라벨이 하나뿐인 위험집합은 제외한다.
로그 손실과 Brier 점수는 확률 예측의 전반적인 정확도를 나타낸다.
단계별 보정 오차는 단계마다 평균 예측확률과 관측 탈락률의 차이의 절댓값을 구해 표본 수로 가중 평균한 값으로, 설계가 단계마다 달라지는 위험 수준을 표현하는지를 보여 준다.
상대 순위에 대해서는 최종 순위와의 Spearman 상관과 상위 5위 판별 AUC를 사용하며, 순위 보조 손실로 학습한 PTR-Former는 순위 출력을, 나머지 모델은 다음 단계 탈락확률 또는 Cox 점수를 사용한다.

\subsubsection{통계적 비교}
두 설계의 차이는 코호트마다 경기 단위로 2,000회 재표집한 쌍대 bootstrap으로 신뢰구간을 구하며~\cite{Efron93}, 이 구간은 학습된 모델을 고정한 상태에서 시험 표본이 달라질 때의 변동을 나타낸다.
코호트별 결과는 DerSimonian--Laird 확률효과 모형으로 통합하고~\cite{DerSimonian86}, REML 추정과 Hartung--Knapp 구간을 민감도 분석으로 함께 확인한다~\cite{Viechtbauer05,HartungKnapp01,SidikJonkman02}.
의미 있는 최소 효과 크기는 AUC 0.005로 둔다.
통합 추정치의 95\% 신뢰구간 하한이 0보다 크고 점추정치가 0.005 이상이면 개선, 95\% 신뢰구간 상한이 0보다 작고 점추정치가 $-0.005$ 이하이면 저하, 90\% 신뢰구간이 $\pm0.005$ 안에 있으면 실질적 동등, 그 밖에는 불확정으로 분류한다.
구조 요인의 다섯 비교는 실험 전에 분석 계획을 고정한 사전 지정 분석이며, 다섯 개의 통합 검정에 Holm 방법을 적용한다~\cite{Holm79}\footnote{다섯 비교는 단계 축과 팀 축의 2$\times$2 요인 실험에서 다른 축의 두 수준에 대해 평균한 단계 축의 주효과, 같은 방식의 팀 축 주효과, 팀 축 주효과의 Spearman 상관 차이, 관계를 섞은 대조 모델 대비 관계 편향의 효과, 서식지 특징을 뺀 입력에서 구역 객체 토큰의 효과이다. 계획을 고정한 뒤 데이터 처리의 누수를 발견하여 고쳤으며, 모든 결과는 고친 자료에서 다시 얻었다. 이 수정은 이전 결과를 확인한 뒤 결정하였다.}.
% 익명 심사본에서는 계획 고정 커밋을 생략하고, 최종본에 추가한다.
생태 hazard 모델, 기존 설계와의 비교, 적용판의 원래 과제 평가, 단일 요인 분해, 단계 안·경기--단계 내부 AUC는 이후에 추가한 탐색적 분석으로 구분하여 보고한다.

\subsubsection{강건성 점검}
신경망 모델은 모든 코호트에서 하나의 난수 시드로 학습하고, \rev{가장 큰 두 코호트에서 시드를 두 개 더 사용하여} 학습 초기화에 따른 변동을 확인한다.
같은 두 코호트에서 학습 횟수의 상한을 두 배로 늘려 상한이 성능을 제한하는지 확인하고, 학습 경기 수를 \rev{네 수준으로} 바꾼 학습 곡선으로 데이터 규모의 영향을 확인한다.
